% Context from chapters-tex/convex-analysis.tex; for mathematical reference, not compilation.
 line $\bar\RR$, with the convention $0\cdot(+\infty)=0$ at the endpoints $t=0,1$.
 
The function $f$ is strictly convex if the inequality in~\eqref{eq-convex-defn} is strict whenever $x,y\in\dom(f)$ are distinct and $0<t<1$.
 
A set $\Om$ is convex if and only if $\iota_{\Om}$ is a convex function.



\begin{figure}[tbp]
\centering
\BookDrawing{tikz/convex-analysis-convex-examples}
\caption{Convexity and strict convexity for functions and sets.}
\label{fig-convex-examples}
\end{figure}




Throughout the chapter, convex functions $f$ are assumed proper, meaning $\dom(f) \neq \emptyset$, and lower semicontinuous (lsc), meaning that for every $x \in \Hh$,
\eq{
	\liminf_{y \rightarrow x} f(y)  \geq f(x).
}
For a convex function, lower semicontinuity is equivalent to the epigraph $\epi(f)$ being closed.
 
We denote by $\Ga_0(\Hh)$ the set of proper convex lsc functions.

                                                                                  
\subsection{First-Order Conditions}

   
\paragraph{Existence of minimizers.}

We first address existence. Lower semicontinuity and coercivity provide a useful sufficient condition.

\begin{prop}
	If $f$ is proper, lower semicontinuous, and coercive (i.e. $f(x)\rightarrow +\infty$ as $\norm{x}\rightarrow +\infty$), then there exists a minimizer $x^\star$ of $f$.
\end{prop}

\begin{proof}
Choose $x_0\in\dom(f)$. The sublevel set $K=\{x:f(x)\leq f(x_0)\}$ is nonempty, closed by lower semicontinuity, and bounded by coercivity. It is therefore compact in finite dimensions. A lower semicontinuous function attains its infimum on $K$, at some $x^\star$. Since $f(x)>f(x_0)\geq f(x^\star)$ outside $K$, this point is a global minimizer.
\end{proof}

This compactness argument is known as the direct method in the calculus of variations.
 
For $f$ in $\Ga_0(\Hh)$, the minimizer set $\argmin f$ is closed and convex, and every local minimizer is global. Strict convexity of $f$ implies that there is at most one minimizer.

   
\paragraph{Subdifferential.}

The subdifferential at $x$ of $f$ is
\eq{
	\partial f(x) \eqdef \enscond{u \in \Hh^*}{ \forall y\in\Hh, f(y) \geq f(x) + \dotp{u}{y-x} },\qquad x\in\dom(f).
}
We set $\partial f(x)=\emptyset$ outside $\dom(f)$. Here $\Hh^*=\RR^N$ denotes the space of dual vectors. The inner product identifies the dual space with $\Hh$, but distinguishing primal and dual variables remains useful. The duality pairing itself is canonical; the identification of a linear functional with a vector depends on the inner product.
 
The subdifferential $\partial f(x)$ consists of the slopes $u$ of supporting affine functions $f(x) + \dotp{u}{z-x}$ that lie below $f$.


\begin{figure}[tbp]
\centering
\BookDrawing{tikz/convex-analysis-subdifferential}
\caption{The subdifferential}
\label{fig-review-convex-analysis-subdifferential}
\end{figure}
At an interior point of $\dom(f)$, the function $f$ is differentiable exactly when its subdifferential consists of the gradient alone:
\eq{
	\partial f(x) = \{ \nabla f(x) \}. 
}
Multiple elements of $\partial f(x)$ indicate distinct supporting slopes of $f$ at $x$.

The subdifferential can also be empty at a boundary point of the effective domain. For example, $f(x)=-\sqrt{1-x^2}$ on $[-1,1]$, extended by $+\infty$ outside this interval, has no subgradient at $x=\pm1$.

Since $\partial f(x) \subset \Hh^*$ is an intersection of half-spaces, it is a closed convex set.
                                                      
 
Thus $\partial f : \Hh \mapsto 2^{\Hh^*}$ is a set-valued operator, also written $\partial f : \Hh \hookrightarrow \Hh^*$.

\begin{rem}[Maximally monotone operator]
The subdifferential $\partial f$ is monotone: setting $U=\partial f$ gives
\eq{
	\foralls (u,v) \in U(x) \times U(y), \quad
		\dotp{y-x}{v-u} \geq 0. 
}
For single-valued maps on the real line, monotonicity means being nondecreasing.
 
Subdifferentials can also be shown to be maximally monotone, meaning that their graphs are not properly contained in the graph of another monotone operator.
 
Conversely, a monotone map need not be a subdifferential; an example is $(x,y) \mapsto (-y,x)$.
  
Many results extend from subdifferentials to arbitrary maximally monotone operators. We restrict attention to subdifferentials here.
\end{rem}

For the absolute value $f(x)=|x|$, write $\partial f(x) = \partial |\cdot|(x)$. Its subdifferential is
\eq{
	\partial |\cdot|(x) = 
	\choice{
		\{-1\} \qifq x < 0, \\
		\{1\} \qifq x > 0, \\
		{[-1,1]} \qifq x=0.
	}
}

\begin{figure}[tbp]
\centering
\BookDrawing{tikz/convex-analysis-subdiff-l1}
\caption{Subdifferential of the absolute value and a piecewise affine convex function.}
\label{fig-subdiff-l1}
\end{figure}







   
\paragraph{First-order conditions.}

The subdifferential gives a necessary and sufficient optimality condition.

\begin{prop}
	$x^\star$ is a minimizer of $f$ if and only if $0 \in \partial f(x^\star)$.
\end{prop}
\begin{proof}
	One has
	\eq{
		x^\star \in \argmin f
		\quad\Leftrightarrow\quad
		\pa{ \forall y, f(x^\star) \leq f(y) + \dotp{0}{x^\star-y} }
		\quad\Leftrightarrow\quad
		0 \in \partial f(x^\star).
	}
\end{proof}

   
\paragraph{Subdifferential calculus.}

Calculus rules simplify subdifferential computations. For a separable function $f(x_1,\ldots,x_K)=\sum_{k=1}^K f_k(x_k)$, the subdifferential is the Cartesian product
\eq{
	\partial f(x_1,\ldots,x_K) = \partial f_1(x_1) \times \ldots \times \partial f_K(x_K).
}
Applying this rule to the $\ell^1$ norm $\norm{x}_1=\sum_{k=1}^N |x_k|$ gives
\eq{
	\partial \norm{\cdot}_1(x) = \prod_{k=1}^N \partial |\cdot|(x_k)
}
a hyperrectangle. With $I = \supp(x)$, the condition $u \in \partial \norm{\cdot}_1(x)$ is equivalent to
\eq{
	u_I = \sign(x_I)
	\qandq
	\norm{u_{I^c}}_\infty \leq 1.
}

A sum rule requires a domain qualification. If one function is finite and continuous at a point in the domain of the other, then for every $x\in\dom(f)\cap\dom(g)$,
\eq{
	\partial (f+g)(x) = \partial f(x) \oplus \partial g(x) = \enscond{u+v}{(u,v) \in \partial f(x) \times \partial g(x)}
}
where $\oplus$ denotes the Minkowski sum. In particular, if $f$ is differentiable at $x$, then
\eq{
	\partial (f+g)(x) = \nabla f(x) + \partial g(x) = \enscond{\nabla f(x) + v }{ v \in \partial g(x) }.
}
Positive scaling gives
\eq{
	\foralls \la >0, \quad
	\partial (\la f)(x) = \la (\partial f(x)). 
}

General compositions need not preserve convexity, so no unrestricted chain rule is available.
 
Composition with a linear map does preserve convexity. For $A \in \RR^{P \times N}$ and $f \in \Ga_0(\RR^P)$, assume $\Im(A)\cap\relint(\dom(f))\neq\emptyset$. Then $f \circ A \in \Ga_0(\RR^N)$ and
\eq{
	\partial (f \circ A)(x) = A^* (\parti